\chapter{Discussion}
\label{ch:discussion}

\section{Answers to the Research Questions}
\label{sec:disc-rqs}
% TODO: answer the guiding question and its four sub-questions from
% Sec.~\ref{sec:research-questions} one at a time, each with the evidence from
% Ch.~\ref{ch:results}. For every evaluated hypothesis state supported / refuted /
% inconclusive -- a refuted hypothesis is a valid finding, not a failure.
% Hypotheses that were registered but not evaluated belong in
% Ch.~\ref{ch:limitations-future-work}, not here.

\section{Reading as the Binding Constraint}
\label{sec:disc-reading-constraint}
% TODO: the interpretive core of the thesis. The attribution result says the
% document-reading stage, not the structuring stage, governs end-to-end quality on
% this task. Discuss what follows for practitioners: effort spent on prompt and
% schema engineering has a low ceiling while the transcript is poor; the cheapest
% large gain is a better reader. Relate to the reported robustness gap of
% open-source parsers on degraded documents.

\section{Measuring Before Optimising}
\label{sec:disc-measurement}
% TODO: the methodological argument. A meaningful share of the initially observed
% error was an artefact of the measurement rather than of the model. Discuss the
% general lesson: on a task where the reference is machine-readable and the input is
% a rendered page, "as printed" and "as stored" diverge systematically, and a naive
% comparison mistakes representation differences for model error. Argue why
% validating the ruler belongs in the method, not in a footnote.

\section{Threats to Validity}
\label{sec:disc-threats}
% TODO: internal validity (leakage guards, deterministic decoding, sealed split);
% construct validity (does the metric measure legal usability? is the reading proxy
% faithful?); external validity (synthetic born-digital corpus vs real scanned
% documents; one language; one document class).
% Limitations proper are collected in Ch.~\ref{ch:limitations-future-work}; keep
% this section about validity threats and cross-reference rather than duplicate.
